


Branch \texttt{redesign} from tag \texttt{combined-v1} (the tag \texttt{release-finish} does not exist). No tagged result and no existing key is changed. New keys use the prefix \texttt{rd.}. Nothing here reads \texttt{data/hardware/}. No synthetic crash pulses are used (no generator passed the realism gate; \texttt{DIAGNOSIS\_SYNTH.md}).

\subsection*{0. What was run before this file (disclosure)}
\begin{itemize}
\item D0 diagnosis (\texttt{rd\_d0.py}, keys \texttt{rd.d0.*}): reads the old design and held-out logs and re-runs cells of the OLD grids with per-alert counters (DIAGNOSIS role). A mechanism sweep of SCT retry spacing was run on the OLD design grid (DIAGNOSIS role).
\item Code checks: the simulator hooks were checked against 750 stored numbers of the old held-out grid, and one smoke run of D1 (3 simulated hours, one cell of the OLD held-out grid: n = 100, p = 0.2, SF8, duty 1\%, b = 30) was run to see that the hook works. The OLD held-out grid is not used to tune anything in this round, so this does not affect the new grids. It is logged in \texttt{DECISIONS.md}.
\item No run on any DEV or FRESH grid exists at the time of writing.
\end{itemize}

\subsection*{1. Integrity rules}
\begin{itemize}
\item \textbf{Three data roles.} DIAGNOSIS: only the existing grids (DESIGN\_A seed 16000, HELD\_A seed 18000) and existing logs. DEVELOPMENT: a new grid (section 3) with new seeds; at most 5 logged iterations per algorithm. FRESH HELD-OUT: a new grid drawn and hashed after the algorithm it tests is frozen and before that algorithm runs on it. The old held-out grid is never used to tune.
\item \textbf{Two fresh grids.} Because the work is done in two sessions, FRESH-A tests D1 and D2 (frozen at the end of session 1) and FRESH-B tests D3 and D4 (frozen in session 2). Each is drawn after the freeze of the algorithms it tests. Each algorithm runs once on its grid. D2 is compared with the frozen D1; D3 with the frozen D2; D4 with R3.
\item \textbf{One Holm family}, at most 10 hypotheses, listed in section 6. A hypothesis is SUPPORTED only if its criterion is met on the fresh grid and its Holm-adjusted p is below 0.05.
\item \textbf{Every result is reported next to the original failure.} A redesign that fails again is a negative result and stays one. The old negative results stay in the negatives register.
\item \textbf{Baselines} are the strongest existing ones: severity table (P1), UA-SCT (P4), fixed 3 retries (P0), quick retries (SCT with 0.3 s spacing, point classes), Tamang-style slots, Khan-style retries (section 4).
\item \textbf{After a freeze}, code and parameters are not changed; findings of reviews after the freeze are recorded without changing results.
\end{itemize}

\subsection*{2. Failures and diagnoses the redesigns answer (from D0; \texttt{docs/03\_results/DIAGNOSIS\_*.md})}
{\scriptsize\raggedright
\begin{longtable}{p{0.317\linewidth}p{0.317\linewidth}p{0.317\linewidth}}
\toprule
Failure & Diagnosis (measured on the old grids) & Redesign \\
\midrule
\endhead
UA-SCT loses to the severity table in steady state (a few cells) and in the single-car 60\% loss case & 5 s retry spacing leaves 4 chances before the 10 s deadline against 10 for the table (single car, 60\% loss, \texttt{t4.b.analytic.*}); in steady state the gain of UA-SCT over the table (held-out mean +0.028 Critical) is no larger than the gain of lowered thresholds (+0.040); on the old design grid SCT with 0.3 s or 1 s spacing beats 5 s spacing by 0.10 to 0.12 in steady state, 1 to 2.5 s spacing is best in bursts (0.3 s collides more, 5 s leaves too few attempts), and the spreading-factor raise adds 0.05 (steady) to 0.09 (burst) & D1 \\
OS-DR loses in 8 of 32 held-out cells & 0.81 to 0.85 of the alerts it suppresses had never reached the base: a vehicle stopped on another vehicle's ACK & D2 \\
GS-CA loses in all 16 burst cells & collisions are the main cause of failure (about half of the failed Critical alerts, a larger share of transmissions); GS-CA put 2.7 times as many frames on the channel as P4+lbt & D3 (addendum) \\
R2 frame size: effect met the threshold in 10 of 16 cells, sign test p = 0.23 & the 16-cell sign test has a power of 0.22 at that effect & D4 uses a paired bootstrap over events \\
\bottomrule
\end{longtable}}

\subsection*{3. Grids}
Cell = (n vehicles, loss p per direction, SF, duty limit, b burst vehicles; b = 1 means no burst). Per cell, simulated hours with new seeds; the arrival process and the class pools are those of \texttt{pa\_sim} (real out-of-fold class pairs, no synthetic pulse); geometry and capture as in \texttt{fin\_g.WORLD}. Paired comparison: the arrival process is the same for every policy (same seed, hour and alert order), so events pair by (hour, alert id).

\begin{itemize}
\item \textbf{DEV grid} (seed base 31000, 100 hours per cell): n in \{80, 180\}; p in \{0.15, 0.35, 0.55\}; SF in \{8, 10\}; duty in \{0.01, 0.002\}; b in \{1, 25\}: 48 cells.
\item \textbf{FRESH-A} (seed base drawn after the freeze; 600 hours per cell): 64 cells = n (2 values) x p (4 values) x SF (2 values) x duty (2 values) x b (1 and one burst size at least 15). The value sets are drawn after the freeze by a seeded generator from the SHA-256 of the freeze hash, within these ranges: n\_low in [55, 95], n\_high in [180, 260]; the four loss levels are 0.15, 0.30, 0.45, 0.60 each shifted by a draw in [-0.03, 0.03] (rounded to 0.01); SF in \{9, 11\}; duty in \{0.005, 0.10\}; burst size in [18, 26]. The grid file and its SHA-256 are committed before the run. The 10\% duty value is the 433 MHz allowance used in the earlier \texttt{duty10} scenario; its regulatory status for India is not verified.
\item FRESH-B is specified in the addendum.
\end{itemize}

\subsection*{4. Policies (exact)}
Metrics are per crash alert whose TRUE class is Critical (deadline 10 s) or Severe (deadline 30 s), measured from the crash instant; "delivered before the deadline" means the base received the alert (an ACK at the vehicle is not needed).

\begin{itemize}
\item \textbf{P1} severity table (policy \texttt{severity}, point classes). \textbf{P0} fixed 3 retries. \textbf{P4} UA-SCT (upper-end classes, fast-path cap phi = 0.1, spacing 5 s / 10 s, SF raise after 2 rounds, the tuned \texttt{s11.tuned.*} values). \textbf{Quick} SCT with point classes and 0.3 s spacing (the earlier \texttt{P2\_quick}). \textbf{Tamang-style} the \texttt{tamang} mode of \texttt{pa\_sim} with the slot cap chosen on the old design grid (K = 8) and the burst size given as the oracle estimate. \textbf{Khan-style} every alert treated as the same class, sent at once, up to 8 retries, delay U(1, 3) s after the ACK window, no spreading-factor raise (re-implemented from the description, an ASSUMPTION).
\item \textbf{D1} on the SCT machinery with UA-SCT's upper-end classes and without the fast-path cap. Beta posterior of the round-failure probability per vehicle (prior mean 0.3, strength 6, forgetting 0.97 per update, updated at every round end from the vehicle's own ACK history; overheard traffic is used only as the load index, see below). Table solved offline by backward induction for each base SF, indexed by (class, failed attempts so far, time left to the class deadline, loss bin, load bin, airtime budget units); actions: retry gap in \{0, 0.25, 0.5, 1, 2, 3.5\} s and SF in \{base, +1, +2\}; objective: probability that the base receives the alert before the class deadline. The link model inside the table: round failure f gives per-direction loss 1 - sqrt(1 - f); SF steps scale per-direction loss by 0.7 (the simulator's SF decay, an ASSUMPTION also used by the table); a quick retry meets the same crowd with a probability that depends on the load bin and decays with the gap (parameters \texttt{rho}, \texttt{g0}). The load index is the number of gateway ACKs the vehicle overheard in the last 10 s (bins 0, 1 to 3, 4 to 9, 10 and more). The event may plan with a share of the remaining hourly allowance (Severe 0.3, Critical 0.5). The table is made monotone by a projection (a retry gap never grows when less time is left or the loss is higher). If no table action fits (deadline or budget), the alert continues on the normal path with exponential backoff and is never stopped except by its try limit (8 Severe, 12 Critical) or an ACK. Firmware cost: the table (about 187 kB at one byte per entry for one SF) plus the two Beta counters.
\item \textbf{D2} on top of D1. In a burst cell the base, when it has heard at least \texttt{thr} distinct vehicles in the last 10 s, replaces per-vehicle ACKs by a cluster ACK listing up to 24 vehicle ids received so far in the cell, at most one per \texttt{tc} seconds, signed with the gateway's Ed25519 key (64-byte signature, 11-byte header, 2 bytes per id); vehicles hold only the gateway public key. A vehicle stops only if its own id is in a verified cluster ACK; if the overhead rate at the vehicle is at least \texttt{thr} and it verifies a cluster ACK that does not list it, it defers its next retry by a class interval (Severe 4 s, Critical 1 s, times a jitter in [0.5, 1] and a scale), at most 3 times per alert; a deferral is never a stop. Outside a burst cell D2 behaves exactly as D1 (tested).
\end{itemize}

\subsection*{5. DEV protocol}
At most 5 logged iterations per algorithm on the DEV grid, written to \texttt{results/tables/rd\_dev\_log.csv} (date, algorithm, parameters, DEV metrics, decision). Tunable: D1: \texttt{rho} scale, \texttt{g0}, prior strength and mean, budget shares. D2: \texttt{thr}, \texttt{tc}, defer scale. Not tunable: the hypotheses, margins, class deadlines, signature length, the baselines. DEV objective, D1: the largest mean per-cell paired gain over P1 among iterations with no DEV cell more than 1 point below P1; D2: the largest mean gain over D1 in burst cells among iterations with no cell more than 2 points below. Iteration 1 uses the parameters stated in section 4; iterations that change a parameter are numbered 2 to 5. The chosen parameters are frozen in \texttt{results/tables/rd\_frozen\_A.json} with SHA-256 of \texttt{rd\_policy.py}, \texttt{world\_sim.py}, \texttt{rd\_eval.py} and the parameter file.

\subsection*{6. Hypotheses and tests (family of at most 10; FRESH-A unless stated)}
Test statistic everywhere: paired bootstrap over events within a cell (2,000 resamples, seed 42), a per-cell one-sided p-value, and an intersection-union rule across cells (the largest per-cell p-value) where a bound must hold in every cell. Pooled statements use the mean over cells of per-cell differences, with events resampled within each cell.

{\scriptsize\raggedright
\begin{longtable}{p{0.237\linewidth}p{0.237\linewidth}p{0.237\linewidth}p{0.237\linewidth}}
\toprule
Id & Statement & Criterion & p-value \\
\midrule
\endhead
H-D1a & No-harm: D1 is not lower than the severity table (P1) by more than 1 point in any FRESH-A cell. & In every cell the paired difference (D1 minus P1, Critical and Severe events pooled, delivered before the own class deadline) has its one-sided 95\% lower bound above -0.01. & largest per-cell p of H0: difference <= -0.01 \\
H-D1b & D1 is at least 5 points above P1 in the cells with loss of 40\% or more. & Mean over those cells of the per-cell difference >= 0.05 with the lower bound of the pooled interval above 0.05. & bootstrap p of H0: pooled difference <= 0.05 \\
H-D1c & D1 never exceeds the duty limit. & No run of D1 in any cell has a trailing-hour airtime above the cell's duty limit (independent scan of the log, \texttt{duty\_exceeded} = 0), and the 10\% cells are included. & 0.999\textasciicircum{}N for N runs with zero violations (H0: a violation rate of 0.001 or more) \\
H-D1s & Strict baseline: against the best baseline of each cell, D1 is not lower by more than 2 points in any cell. & The best of \{P1, P4, P0, Quick, Tamang-style, Khan-style\} is taken per cell on the same fresh data (this favours the baselines); the lower bound of D1 minus that baseline is above -0.02 in every cell. & largest per-cell p of H0: difference <= -0.02 \\
H-D2a & No-harm: in no FRESH-A cell does D1 + D2 lose more than 2 points of final delivery against D1 alone. & Final delivery = the base received the alert at any time before the end of the horizon, Critical and Severe pooled; lower bound of the paired difference above -0.02 in every cell. & largest per-cell p of H0: difference <= -0.02 \\
H-D2b & In burst cells D2 raises delivery within 10 s by at least 5 points over D1 alone. & Critical events, the base received the alert within 10 s of the crash; mean over burst cells of the per-cell difference >= 0.05 with the lower bound above 0.05. & bootstrap p of H0: pooled difference <= 0.05 \\
H-D2c & In the bounded model, a node holding only a vehicle key cannot make another vehicle stop retrying. & Explicit-state check (\texttt{rd\_d2\_model.py}): two or three vehicles, one gateway, an attacker with at most 3 actions, at most 2 losses, replays of older frames and tampering; no vehicle stops without a verified listing, defers are bounded; the limit case of a stolen gateway key is outside the claim and reported. & 0 if the check passes, 1 otherwise \\
H-D3, H-D4a, H-D4b & Specified in the addendum before their development (the family is closed at 10). &  &  \\
\bottomrule
\end{longtable}}

\textbf{Kill criteria.} D1: H-D1a or H-D1b fails. D2: any of H-D2a, H-D2b, H-D2c fails. A killed redesign is reported as a negative result next to its original failure.

\textbf{Wording.} Plain language. "v2 designed from a diagnosis of v1 failure, tested once on a fresh held-out grid" only for supported redesigns. Absence of prior work is stated only as "no study found in the sources checked". Every effect carries a 95\% interval. Every table caption points to \texttt{docs/03\_results/ASSUMPTIONS.md}.

\subsection*{7. Code verification (CV; \texttt{docs/03\_results/CODE\_VERIFICATION.md})}
Before any DEV or FRESH run of a module: unit tests with hand-computed cases; property tests (hypothesis); run-time invariants in the simulator (\texttt{invariants=True} aborts the run); a differential test of the simulator policy against the firmware-path implementation (D6, session 2) on 10,000 random traces; mutation testing with at least 20 mutants of the policy and security modules (below 80\% killed means the tests are extended before evaluation); an independent review in a fresh context; reproduction of three earlier numbers (a UA-SCT cell, an OS-DR cell, the R3 result) with the changed simulator.

\subsection*{8. Cross-simulator check (X1; written now, run in session 2)}
ns-3 with the signetlabdei \texttt{lorawan} module (build recorded in \texttt{docs/03\_results/XSIM\_BUILD.md}). Calibration: the same baseline scenarios in both simulators (single gateway; 1, 10, 50, 100 nodes; SF7 and SF9; confirmed and unconfirmed uplinks). Agreement criterion, fixed now: the absolute difference in packet delivery ratio is at most 5 points in every baseline cell. If it is not met, the cross-simulator verdicts are reported as inconclusive. A D-hypothesis whose verdict differs between the simulators is reported as simulator-dependent and cannot be in the abstract.

\subsection*{9. Addendum A (2026-11-14, written after the independent code review and before the freeze of D1 and D2; no DEV or FRESH-A result changes any hypothesis)}
Findings of the review are in \texttt{docs/03\_results/CODE\_VERIFICATION.md} section 6. The review found no blocker and no use of held-out data. Changes and disclosures made before the freeze:

\begin{itemize}
\item \textbf{Statistic made more conservative.} Events of one simulated hour share a channel episode (in a burst, 18 to 26 Critical alerts), so an event bootstrap overstates the evidence. Every per-cell and pooled test is now computed twice, with events resampled (section 6) and with simulated hours resampled within each cell (cluster bootstrap); the verdict uses the larger p-value and the smaller lower bound. Both are reported. Margins, directions and the intersection-union rule are unchanged.
\item \textbf{Pooled criterion form.} H-D1b and H-D2b are met when the point estimate is at least 0.05 and the one-sided p-value of H0 (difference at most 0.05) is below 0.05. The reported 95\% interval is two-sided and therefore slightly stricter.
\item \textbf{Bounds of the H-D2c model} (section 6 said two or three vehicles, at most 3 attacker actions, at most 2 losses): run with 2 vehicles at those bounds, and with 3 vehicles at attacker actions at most 2, losses at most 1 and one send per vehicle (the larger state space does not fit this representation). Sends are bounded by 2 (2 vehicles) and defers by 3 (as in the simulator). The model abstracts the signature check, so cryptography is tested by unit tests, not by the model; the defer bound holds largely by construction; the freshness window cannot reject a frame inside the model horizon; the attacker's own alerts are valid only under its own vehicle key (R4 keys). A stolen gateway key is reported as a limit case.
\item \textbf{Code fixes before the freeze} (behaviour now matches section 4): a demoted D1 alert keeps the class try limit (8 Severe, 12 Critical); a defer cannot push a retry past the class deadline minus one round at the base SF; a demoted alert returns to the base SF; delays are kept in 64-bit floats; D2 requires thr of at least 2. After these fixes the chosen parameters were re-run on the DEV grid as a check, not as an iteration (\texttt{results/tables/rd\_dev\_checks.csv}); nothing was re-chosen.
\item \textbf{Disclosed limitations, not changed.} (a) In the simulator a retry round is a 14-byte flash followed by the alert (the SCT machinery), while the D1 table plans one 97-byte transmission per attempt: the table underestimates the time of a round and ignores that a flash reception already counts as reception. (b) The Beta mean already includes the failures of the current event and the table adds them again, so the policy grows more pessimistic as an event ages. (c) The table caps are NB\_MAX = 8 budget units and J\_MAX = 8 failed attempts (the Critical try limit is 12). (d) The cluster ACK lists the 24 most recently received vehicles, so in a burst of 25 or 26 vehicles the earliest receivers drop off until they retry. (e) The simulator does not call \texttt{ClusterAckVerifier} or \texttt{NodeState}; their equivalence to the simulator logic is the subject of the D6 differential test (session 2). (f) DEV and FRESH-A share the same finite pool of real out-of-fold class pairs and early-exit decision times; only arrivals, channel draws and geometry are new, so a fresh grid is fresh in those, not in the class pool. (g) The initial values of \texttt{rho}, \texttt{g0}, \texttt{thr}, \texttt{tc}, the loss and load bin edges, the table step sizes and the defer intervals are design values with no data origin; iteration 1 uses them as written and later iterations are in \texttt{rd\_dev\_log.csv}.
\item \textbf{Freeze coverage extended:} the freeze records SHA-256 hashes of \texttt{rd\_policy.py}, \texttt{world\_sim.py}, \texttt{rd\_eval.py}, \texttt{rd\_d2\_model.py}, \texttt{rd\_common.py}, \texttt{rd\_fresh\_a.py}, \texttt{rd\_dev.py} and of the modules that define the baselines and the physics (\texttt{fin\_g.py}, \texttt{fin\_sim.py}, \texttt{pa\_sim.py}, \texttt{s11\_sct.py}, \texttt{sct\_sim.py}, \texttt{lora.py}, \texttt{capacity.py}, \texttt{packet.py}, \texttt{replay.py}).
\item \textbf{DEV outcome known at the time of this addendum:} D1 used 3 of 5 iterations, D2 used 5 of 5 (\texttt{rd\_dev\_log.csv}). The best D2 iteration gained about 4.5 points in the DEV burst cells against the 5 points that H-D2b requires, so H-D2b may fail on FRESH-A; it is run as preregistered.
\end{itemize}


\subsection*{10. Addendum B (2026-11-14, session 2; written before any D3 or D4 development run, before any firmware-path or ns-3 run of D1, and before FRESH-B is drawn)}

Nothing in sections 1 to 9 is changed. Session 2 runs after the FRESH-A result (D1 supported on H-D1a, H-D1b, H-D1c, H-D1s; D2: H-D2a and H-D2c supported, H-D2b not supported, so D2 is a negative result). Code and parameters of D1 and D2 stay frozen (\texttt{results/tables/rd\_frozen\_A.json}); nothing here re-tunes them.

\subsubsection*{10.1 Contradiction between the task text and the repository (D4), and how it is handled}
The task text describes R3's provisional alert as a 97-byte frame, and asks D4 to make extra airtime fall from 38\% to below 25\% by sending a 24-byte provisional frame. In this repository R3's provisional alert is a \textbf{14-byte flash}, the confirm is the 97-byte alert, and a retraction is a 14-byte authenticated CANCEL sent three times (\texttt{ar\_r3.py}, V11 section 4; the 38\% is \texttt{ar.r3.a1.air\_ratio\_p0} = 0.380). A 24-byte provisional frame is longer than a 14-byte flash (SF7: 61.7 ms against 46.3 ms), so it cannot lower the extra airtime of R3 as registered. The conservative handling (logged in DECISIONS.md): D4 is registered as the task describes it; H-D4a is tested against 25\% exactly as worded; the comparator for H-D4b is named as R3 with the 97-byte provisional frame (the task wording) and also as R3 as registered. Nothing here was computed from FRESH-B before this section was committed; the arithmetic above uses only airtime values.

\subsubsection*{10.2 D3: gated slotted retries with an overheard burst-size estimate (replaces GS-CA)}
Mechanism (D0, \texttt{DIAGNOSIS\_GSCA.md}): 0.54 of the failed Critical alerts of GS-CA failed mainly through a collision or a deaf gateway, and slotting with a fixed K put 2.7 times as many frames on the channel as P4+lbt. D3 is therefore kept.

Design, on top of the frozen D1 + D2 (policy \texttt{D1D2}; D3 uses the frozen parameters of FRESH-A):
\begin{itemize}
\item \textbf{No oracle.} Vehicle i keeps N\_hat(i), the number of distinct other vehicles whose alert or flash frames it decoded in the last 10 s (node-to-node link as in \texttt{world\_sim.node\_link}), plus itself. The initial transmission of every alert is exactly as D1 (a vehicle cannot know a burst is under way before it hears one); slots apply to \textbf{retries only} (attempt index j >= 1) of Critical alerts whose N\_hat >= \texttt{thr3}.
\item \textbf{Slot.} S = clip(ceil(N\_hat / \texttt{m}), 2, S\_max), with S\_max limited so that the epoch (S slots) is at most 5 s (slot = time on air of the alert at the chosen SF + 0.02 s guard). Slot index = SHA-256(GPS second of the epoch start, vehicle id) mod S. The retry is sent at the start of the chosen slot in the earliest epoch whose slot start is not earlier than the D1 gap after the previous round. The own clock offset shifts the transmission by at most 5 ms (as GS-CA). If the slot would end after the class deadline, D1's own gap is used (no slotting). A slot only delays a retry; the try limit and the ACK rule of D1 and D2 are untouched, so slotting never stops an alert.
\item \textbf{Spreading factor.} If \texttt{alt} is on, retry j uses the SF chosen by D1 on odd j and that SF plus one on even j (at most 12); frames of different spreading factors are treated as quasi-orthogonal by the simulator (the simulator's assumption, \texttt{ASSUMPTIONS.md}).
\item Tunable on the DEV grid (at most 5 iterations): \texttt{thr3} (default 3), \texttt{m} (default 2), \texttt{alt} (default on). Not tunable: the epoch cap of 5 s, the guard, the hypothesis, the margins.
\end{itemize}

DEV grid for D3: seed base 51000, 100 hours per cell, the 24 burst cells of the DEV grid of section 3 (n in \{80, 180\}, p in \{0.15, 0.35, 0.55\}, SF in \{8, 10\}, duty in \{0.01, 0.002\}, b = 25). Objective: the largest mean paired gain of D3 over D1+D2 in Critical within 10 s among iterations with no DEV cell more than 2 points below D1+D2 in delivery before the own deadline. FRESH-B for D3: 32 burst cells drawn after the freeze (section 10.5), 300 hours per cell, seeds new. Non-burst cells are not run: D3 is inactive without a burst by construction (tested), and the DEV grid has no non-burst cell for it.

\textbf{H-D3.} In the FRESH-B burst cells D3 raises Critical delivery within 10 s by at least 5 points over D1+D2 (mean over cells of the per-cell paired difference >= 0.05 and the one-sided p of H0: difference <= 0.05 below 0.05, the more conservative of the event and hour-cluster bootstraps), and in no cell is the lower 95\% bound of the paired difference in delivery before the own class deadline (Critical and Severe pooled) below -0.02. Kill: either part fails. Original failure kept beside it: GS-CA lost in all 16 burst cells of the old held-out grid.

\subsubsection*{10.3 D4: short provisional frame in the two-stage alert}
Design as in the task: the provisional alert is the 24-byte frame with an 8-byte tag (fields: class, coarse position, freshness, tag); the confirm is the full 97-byte frame; a retraction is the 14-byte authenticated CANCEL, three transmissions (as R3).
\begin{itemize}
\item \textbf{H-D4a (extra airtime).} On the de-duplicated test events of the 1\% false-alarm operating point (the R3 primary operating point; events and decisions are those of \texttt{s6\_early}, nothing is refitted), the extra transmit airtime of D4 over the single-stage alert (provisional frame, confirm and CANCEL counted; tries counted as in \texttt{ar\_r3.py}) is below 0.25. Grid: 48 cells = 6 spreading factors (7 to 12) x 8 loss levels per direction; the 8 loss levels are drawn after the freeze (section 10.5) within [0, 0.7]. ANALYTICAL on the stored test events; paired bootstrap over events (2,000 resamples, seed 42) per cell; criterion: the upper 95\% bound of the ratio is below 0.25 in every cell; p = the largest per-cell p of H0: ratio >= 0.25.
\item \textbf{H-D4b (arrival time and burst delivery).} The provisional alert's median arrival time is not later than R3's by more than 0.1 s, and Critical delivery within 10 s in burst cells is at least 3 points above R3 with the 97-byte provisional frame (paired bootstrap over events). It needs a two-stage alert inside \texttt{world\_sim} with per-alert frame lengths, which the simulator does not have. \textbf{It is run only if H-D4a is met} (D4 is killed if either part fails, so a failed H-D4a ends D4); otherwise it is reported NOT\_RUN with that reason and enters the Holm family at p = 1.
\item Exploratory, not a hypothesis: extra airtime of the 14-byte flash (R3), the 24-byte frame (D4) and the 97-byte frame as the provisional alert, by spreading factor and loss level, so that the reason for a failure is visible.
\end{itemize}

\subsubsection*{10.4 X1.3 and X3: cross-simulator check and added baselines}
\begin{itemize}
\item \textbf{Scenario (S-X), identical in both simulators.} One gateway at the centre of a disc of 3,000 m; n vehicles placed uniformly; heartbeats of 16 bytes every 600 s; every crash event is Critical (10 s deadline), decided 0.3 s after the crash; Poisson crash arrivals at 0.1 per vehicle-hour, plus, in burst cells, b vehicles with a Critical event at a uniform time in [600 s, 610 s]; extra loss p per direction (uplink and downlink) on top of the physical channel of \texttt{ASSUMPTIONS.md}; duty limit as a trailing-hour airtime budget; spreading factor 9 (ns-3's ALOHA region has no valid data rate for SF11, so the SF11 cells of FRESH-A are not expressible). Cells: the 32 SF9 cells of FRESH-A (grid file \texttt{rd\_fresh\_A\_grid.json}), 100 hours per cell and policy at most, new seeds. These cells were already used once for the custom-simulator verdicts; the ns-3 run uses the frozen D1 table and parameters and re-tunes nothing, so it is a cross-check, not a second held-out test, and its outcome cannot change the FRESH-A verdicts.
\item \textbf{Policies in ns-3.} P1 (severity table, no flash), quick retries (0.3 s spacing, spreading factor raise after 2 rounds), D1 (frozen table of \texttt{rd\_policy.py} through the C firmware-path code, Beta loss estimate, load index from overheard gateway ACKs), and standard LoRaWAN confirmed uplink (the module's class A confirmed uplink with NbTrans = 8, ACK in RX1 or RX2, ACK-timeout retransmission by the module; ADR off, and ADR on if the module's ADR component runs in the scenario). The ns-3 versions of P1 and quick retries are simplified re-implementations from the policy descriptions (they have no flash frame); the custom-simulator versions are the frozen ones with a flash. These differences are stated in the results; each simulator's D1 verdict is relative to its own P1.
\item \textbf{Custom simulator on S-X.} The frozen policies P1, P4, P0, quick, tamang, khan, D1 run on S-X in the custom simulator too (events Critical only, the same cell parameters, new seeds), because the FRESH-A numbers include Severe events and random class mixes and cannot be compared one to one. D1 and every baseline are the frozen code.
\item \textbf{Agreement rule (stated now).} For each of H-D1a (no cell where D1 is lower than P1 by more than 1 point, lower bound), H-D1b (gain of at least 5 points over P1 pooled over cells with loss of 40\% or more), H-D1c (duty limit never exceeded) the verdict is computed in each simulator on S-X by the V12 criteria; the verdict is "same" if both give the same SUPPORTED or NOT SUPPORTED outcome. The calibration criterion (X1.2) was NOT MET (one cell 5.4 points out), so every cross-simulator verdict is labelled with that: an agreement is reported as agreement under a calibration that did not reach its own criterion. A hypothesis whose verdict differs is reported as simulator-dependent and does not go into the abstract. D2 is not run in ns-3 (it is already a negative result).
\item \textbf{X3 baselines.} The strongest baselines of V12 section 4 stay. Added: (a) standard LoRaWAN confirmed uplink as above (ns-3); in the custom simulator the Khan-style policy (up to 8 retries after an ACK timeout drawn from U(1, 3) s, no spreading-factor change) is the closest existing policy to NbTrans = 8 with ADR off, which is stated and no new custom policy is added; (b) the closest adaptive scheme of X2 (\texttt{docs/06\_references/CLOSEST\_PRIOR\_WORK.md}): that file says which of the five papers could be implemented from its text; one that cannot is reported as not implemented, not approximated silently. Per-cell best-baseline testing (H-D1s) already uses the six baselines on FRESH-A; on S-X the best of all baselines of each simulator is used.
\end{itemize}

\subsubsection*{10.5 FRESH-B, freeze and procedure}
\begin{itemize}
\item Freeze: \texttt{results/tables/rd\_frozen\_B.json} records the DEV iteration chosen for D3 and the SHA-256 of every code file listed in section 9 item 6 plus the D3 and D4 modules and the D6 firmware-path files. The freeze is committed before the grid is drawn.
\item Draw: the grid (D3 burst cells: n\_low in [55, 95], n\_high in [180, 260], loss levels 0.15, 0.30, 0.45, 0.60 each shifted by a draw in [-0.03, 0.03], SF in \{9, 11\}, duty in \{0.005, 0.10\}, burst size in [18, 26]; seeds from the hash; D4 loss levels) is generated from the SHA-256 of the freeze file by a seeded generator and committed with its hash. The FRESH-A draw used a different freeze hash, so values differ.
\item Run once: D3 and D4a each once, a marker file written before the run.
\item \textbf{One Holm family} of ten: H-D1a, H-D1b, H-D1c, H-D1s, H-D2a, H-D2b, H-D2c, H-D3, H-D4a, H-D4b. Final Holm p-values are computed after FRESH-B and replace the provisional ones of session 1. A hypothesis not run enters at p = 1.
\item \textbf{Code verification} (section 7) for every new module (D3, the D4 analysis, the C firmware path, the ns-3 application): unit tests with hand-computed cases, property tests, run-time invariants, the differential test of the simulator policy against the firmware-path implementation on 10,000 random traces (target 0 mismatches), mutation testing (the hand-written mutant set; mutmut does not run on Windows), an independent review in a fresh context, and the known-answer reproduction. Findings after the freeze are recorded, not fixed.
\item \textbf{Firmware in an emulator (D6).} Host build (gcc, WSL) of the decision path: alert framing, HMAC-SHA256 tag, freshness, D1 lookup and Beta update, ACK and cluster-ACK verification (Ed25519 verification with TweetNaCl, public domain). Espressif QEMU and the ESP-IDF toolchain are not installed here; if they cannot be installed within the session the emulator and the ESP32 flash and RAM sizes are NOT\_RUN and are not estimated. Emulator timing, if it were run, is labelled EMULATED and never reported as device timing.
\end{itemize}

\subsubsection*{10.6 Clarifications and disclosures after the independent review (written before the FRESH-B freeze; no FRESH-B result exists)}
The D3 DEV runs (4 of 5 iterations, \texttt{rd\_dev\_log.csv}) were made before the independent review of the session 2 code. The review (\texttt{CODE\_VERIFICATION.md} section 9) found no use of held-out data. Where the code differs from sections 10.2 to 10.4 the code is not changed after DEV (a change would need a new DEV round and only 1 iteration is left); the difference is disclosed here instead, and H-D3 is evaluated for the code as it is.
\begin{itemize}
\item \textbf{Simulator files changed after the FRESH-A freeze.} \texttt{world\_sim.py} and \texttt{rd\_eval.py} (in the FRESH-A hash list) gained the D3 hook, the optional decision trace and the D3 policy branch. The D1 and D1+D2 paths are guarded by \texttt{d3} and \texttt{rtrace} and use a separate random stream for the D3 clock offsets. Equivalence check: D1 and D1+D2 reproduce the stored FRESH-A per-event delays bit for bit (4 cells x 8 hours, \texttt{tests/redesign/test\_rd\_regression\_a.py}; earlier 3 cells x 12 hours). The FRESH-A hashes in \texttt{rd\_frozen\_A.json} no longer match these two files; the FRESH-A verdicts are unchanged.
\item \textbf{thr3 counts other vehicles.} Section 10.2 says N\_hat includes the vehicle itself and slots are on when N\_hat >= \texttt{thr3}. The code switches slots on when the number of OTHER vehicles heard is >= \texttt{thr3}, and passes N\_hat = others + 1 to the slot count. \texttt{thr3} was chosen on DEV under this reading. With the literal reading the same \texttt{thr3} would mean one more vehicle was needed.
\item \textbf{Overhearing model.} A vehicle records a frame only while it has a pending alert and is not itself transmitting or in its own listen window; frames heard before its own crash are not counted; no collisions are modelled at the vehicle receivers (so N\_hat may be too high when frames collide, and too low for a late crasher in a burst). Not in \texttt{ASSUMPTIONS.md} before; added.
\item \textbf{The slot covers the initial frame of a round.} In the simulator a retry round is a 14-byte flash followed by the 97-byte alert; the slot time is applied to the flash and the alert follows one listen window later, unslotted. Slot length and the deadline test use the alert's time on air.
\item \textbf{Epoch cap and S\_MAX.} At least two slots are kept even when two slots are longer than 5 s (an SF12 slot is about 3.96 s), so the 5 s epoch cap does not hold at SF12 (reached through \texttt{alt} in SF11 cells). S\_MAX = 12 was written with the module, before any DEV run; it binds only at short slots.
\item **\texttt{alt} and the budget plan.** The alternated spreading factor is not part of D1's planned action, so the per-event planning share of section 4 does not cover it; the trailing-hour duty limit still applies and is checked (H-D1c style scan in the run).
\item \textbf{Code changes made before the freeze because of the review (not tuning):} H-D4b wording in the analysis (NOT\_RUN either because H-D4a failed or because the required simulator feature does not exist); an analyse-only entry point and a pre-flight check before the single run; the freeze falls back to the last D3 iteration, with the DEV floor result recorded, if no iteration met the floor; the freeze hash list now includes \texttt{build.sh}, \texttt{size.sh}, \texttt{results.py}, \texttt{ar\_r1.py}, the class-pool pickle and the design JSON; C hardening (tag check requires the exact frame length, the freshness ring rejects R outside 1 to 512, the harness no longer reads uninitialised tokens); the differential generator uses seeded Ed25519 keys (the earlier runs used unseeded keys). None of these touches a hypothesis, a margin or a parameter.
\item \textbf{Known differences kept:} the C cluster-ACK check rejects frames longer than 255 bytes as malformed before the signature check (Python checks the signature before the length); LoRa frames cannot be longer. The differential test does not cover Python-only behaviour outside the C path.
\end{itemize}

\subsection*{11. Addendum C (2026-11-14, after the redesign-v1 tag; written before any run of the items below)}

The second Session 2 prompt asked for items A1 to A5 before the rest of the Session 2 list. The rest (D3, D4, D6, D7, X3, X4, the reviewer pack, the decision table, tag \texttt{redesign-v1}) had already been done when that prompt arrived, so the order was reversed. This addendum covers A1 to A4 and A7. \textbf{Nothing here changes a committed result, a frozen file of FRESH-A or FRESH-B, or a verdict.} Everything added here is labelled \textbf{ADDED-AFTER-SESSION-1} in every table that shows it, and none of it enters the Holm family of ten (the family is closed; these are baseline additions and a calibration, reported with unadjusted p-values or intervals).

\subsubsection*{11.1 A1: Nessa-style baseline}
\textbf{Source level.} The CNS 2020 paper (Nessa, Hussain, Fernando) was not opened. The scheme is taken from a secondary description, section 3.1 and 5.3 of a Chalmers master's thesis (odr.chalmers.se, opened as text on 2026-11-14), which implements it from the paper. The implementation here is an interpretation of that description, not a reproduction of the paper, and is labelled so.
\textbf{Scheme as implemented (assumptions marked).} After an alert transmission with no ACK the vehicle waits an ACK timeout drawn uniformly in [Lb(i), Ub(i)] of stage i, with Ub(0) = T0 = k x (time on air of the 97-byte alert at the cell's base spreading factor) and Lb(0) = T0 / 2 (ASSUMPTION: the description gives no Lb(0)); Lb(i) = Ub(i-1) and Ub(i) = 2 Ub(i-1) for later stages (doubling, as in the description). It then retransmits twice at the same spreading factor, the two copies separated by a random "arbitrary spacing" drawn uniformly in [0, A]; after the pair it raises the spreading factor by one and starts the next stage. The timeout is measured from the end of the transmission, so the extra wait added after the listen window is max(0, timeout - window). The scheme has no severity classes and no deadline: every alert is treated alike (as the Khan-style baseline does), at most 12 tries, no flash frame, no hold, no duty-cycle reservation (the trailing-hour duty limit of the simulator still applies).
\textbf{Free parameters, chosen on the DEVELOPMENT grid only} (the DEV grid of section 3, seed 31000, 100 hours per cell, 48 cells; at most 5 logged iterations in \texttt{rd\_dev\_log.csv} with algorithm "Nessa"): k in \{1, 2, 4, 8\} and A in \{0.3, 1, 2\} s. Iteration 1 is k = 2, A = 1 s. Objective: the largest mean delivery before the own class deadline over the DEV cells. The parameters are frozen in \texttt{results/tables/rd\_frozen\_C.json} with code hashes before the baseline touches any fresh grid.
\textbf{Runs.} (i) The Nessa-style policy runs once on the FRESH-A grid (frozen file \texttt{rd\_fresh\_A\_grid.json}, 600 hours per cell, same seeds, so events pair with the stored runs of D1 and the other baselines): this adds a baseline and tunes nothing in D1. (ii) The same scheme with the same frozen parameters as an ns-3 application on scenario S-X.
\textbf{Reporting.} D1 minus Nessa-style per cell (paired, the more conservative of the event and hour-cluster bootstraps) with 95\% intervals, the pooled difference, and the share of cells in which D1 is better by at least 2 points (point estimate) and by at least 2 points with the lower bound above 0.02. H-D1s is re-run with seven baselines (the six of section 4 plus Nessa-style): the best of the seven per cell, D1 not lower than it by more than 2 points in any cell (lower bound above -0.02), labelled ADDED-AFTER-SESSION-1. \textbf{Reading rule fixed now:} D1 is called "better than Nessa-style" only if the pooled difference is at least 0.02 with a lower bound above 0 AND at least half of the cells show D1 better by 2 points; otherwise Nessa-style "matches or beats" D1 and D1 is not claimed as an improvement over prior work.

\subsubsection*{11.2 A2: the baseline reading}
Every D1 table, the abstract and the claims state that the severity table is a weak baseline in this simulator (it waits 5 s before the initial Critical send and retries without spacing) and that most of D1's gain over it is also available from quick retries (D1 minus quick retries is about +0.035 on FRESH-A). The abstract and claims use D1 against the best baseline per cell, not against the table.

\subsubsection*{11.3 A3: calibration under two sensitivity settings}
The criterion is unchanged (absolute PDR difference at most 5 points in every baseline cell) and is not relaxed; no custom-simulator constant is changed to pass it. Two settings are fixed now and both are reported: \textbf{(a)} both simulators at the custom simulator's gateway sensitivity (-123, -126, -129, -132, -134.5, -137 dBm for SF7 to SF12; the setting of \texttt{XSIM\_CALIBRATION.md}); \textbf{(b)} both at ns-3's shipped gateway sensitivity (-130, -132.5, -135, -137.5, -140, -142.5 dBm, 7 dB better at SF7). In (b) the ns-3 end-device table is set to the same values (the shipped end-device table is 6 dB worse than the shipped gateway table and the custom simulator has one table for both directions), a stated deviation from "shipped". Both with the local ADR backoff disabled, same scenarios as section 8, 60 hours or runs per cell. The runtime switch in ns-3 is an offset added to the compiled table (a local change recorded in \texttt{XSIM\_BUILD.md}). The failed original criterion (15 of 16 cells within 5 points, one at 5.4) stays in the paper as a limitation. Setting (a) is the setting of all X1.3 verdicts; setting (b) is used for X1.3 only if time allows, as a sensitivity check.

\subsubsection*{11.4 A4: X1.3}
D1, quick retries, the severity table, standard LoRaWAN confirmed uplink (NbTrans 8, ADR off; ADR on NOT\_RUN) and the Nessa-style baseline run as ns-3 applications on every FRESH-A scenario ns-3 can express (the 32 SF9 cells; SF11 is not expressible), 100 hours per cell, setting (a). H-D1a, H-D1b, H-D1c and H-D1s are evaluated as in section 10.4, with the Nessa-style policy added to the baselines of H-D1s. \textbf{D1 goes to the main paper only if H-D1b and H-D1s hold in ns-3 as well; otherwise it is reported as simulator-dependent.} A5 (a D2-v3 variant with 4-byte per-vehicle tags) is NOT\_RUN.

\subsubsection*{11.5 A7: disclosures carried into the paper}
(1) The fresh grids share the finite pool of real class pairs and early-exit decision times with development; freshness holds for arrivals, channel and geometry. (2) The H-D2c bounded model is exhaustive for 2 vehicles at the preregistered bounds; for 3 vehicles the bounds were reduced (attacker actions at most 2, losses at most 1, one send per vehicle). (3) The test statistic is the more conservative of the event bootstrap and the hour-clustered bootstrap. (4) Three older tests failed with numpy-2 attribute errors (\texttt{np.trapz} removed); a compatibility shim (\texttt{trapezoid} when present, else \texttt{trapz}) replaces the three calls in \texttt{fin\_synth.py} and \texttt{ph\_gen.py} without changing any value, and the three tests are re-run.

